\chapter{Conclusiones y Trabajo Futuro}
\label{ch:conclusiones}

\section{Conclusiones}

Este trabajo se propuso diseñar e implementar un sistema capaz de extraer conocimiento estructurado
de reuniones técnicas, gestionar su evolución en el tiempo mediante un grafo, e incorporar la
revisión humana y la evaluación calibrada como requisitos de calidad. El resultado, \textbf{Seshat},
es un sistema completo y desplegable que cumple ese objetivo general. Las conclusiones se organizan
en torno a los cinco objetivos específicos planteados en la introducción.

El \textbf{flujo multi-agente de extracción} (objetivo~1) se materializó en dos pasadas: una de
identificación, con un agente especializado por tipo de concepto, y otra de resolución, que infiere
las relaciones semánticas de cada nodo nuevo con el conocimiento preexistente recuperado por
similitud. La inversión de la premisa habitual de la \ac{RAG} —usar la recuperación no para
responder preguntas, sino para construir y mantener el propio grafo— resultó ser la decisión
conceptual que articula todo el sistema.

La \textbf{revisión humana asistida por confianza} (objetivo~2) se implementó como una compuerta
entre ambas pasadas: cada nodo recibe una puntuación heurística que decide si se aprueba de forma
automática o se deriva a un revisor, cuyo umbral se calibra —en vez de fijarse a mano— para
maximizar la cobertura sin comprometer la precisión de lo aprobado. El sistema puede así operar de
forma plenamente automática cuando la confianza lo permite, o concentrar la atención humana en los
nodos de menor certeza.

El \textbf{marco de evaluación} (objetivo~3) es, quizá, la contribución más transferible del
trabajo: cinco arneses independientes miden cada etapa no determinista del \textit{pipeline} contra
un corpus etiquetado, con \textit{scorers} deterministas que garantizan la reproducibilidad. Su
veredicto no es meramente informativo, sino que se materializa en un \textit{gate} que impide el
arranque del sistema si la calidad medida cae por debajo de los umbrales fijados
(capítulo~\ref{ch:evaluacion}). En la última ejecución, las cinco etapas superan sus objetivos de
calidad.

La \textbf{API REST y la interfaz web} (objetivo~4) cubren el ciclo de vida completo de un trabajo
—envío, transcripción, extracción, revisión y consulta del grafo— sin exigir acceso directo a la
base de datos: la \ac{API} es el único punto de entrada, y la interfaz una capa de conveniencia que
delega íntegramente en ella. Por último, el \textbf{despliegue contenerizado con observabilidad}
(objetivo~5) integra la medición del consumo de \textit{tokens}, la latencia por etapa y las trazas
de ejecución en \href{https://mlflow.org}{MLflow}, haciendo el rendimiento y el coste comparables
entre versiones.

En conjunto, \textbf{Seshat} demuestra que un sistema basado en \ac{LLM} puede alcanzar una fiabilidad
apta para producción no por la infalibilidad de sus modelos, sino por la disciplina de ingeniería
que los rodea: el anclaje de cada afirmación a su evidencia, la superposición de salvaguardas contra
la alucinación, y —sobre todo— la conversión de la calidad en una magnitud medible, reproducible y
exigible como condición de despliegue. Las limitaciones conocidas y las líneas de continuación se
detallan a continuación.

\section{Trabajo futuro}
\label{sec:trabajo-futuro}

El desarrollo dejó abiertas varias líneas de continuación —incluidas las capacidades excluidas del
\ac{MVP} (sección~\ref{sec:alcance})—, que se organizan aquí en cuatro categorías según la razón por
la que permanecen pendientes: el endurecimiento que exigiría un despliegue de producción, las
extensiones del núcleo generativo aplazadas por tiempo o por depender de datos reales, la
profundización del grafo de conocimiento —trabajo de modelado de datos más que de \ac{IAG}—, y la
integración con sistemas externos y nuevas fuentes de entrada.

\subsection{Endurecimiento para producción}

Capacidades que un despliegue real exigiría, pero que exceden lo que un prototipo de investigación
necesita demostrar.

\begin{itemize}
  \item \textbf{Despliegue en infraestructura cloud.} El sistema corre en local contra LocalStack;
        llevarlo a una nube real, con su infraestructura definida como código, queda pendiente.

  \item \textbf{Autenticación con proveedor de identidad externo (\textit{SSO}).} El modelo de
        claves de \ac{API} basta para la base de usuarios prevista; un IdP externo —por ejemplo con
        JWT— sería el paso natural al crecer o al necesitar inicio de sesión único.

  \item \textbf{Cola de tareas durable y escalado horizontal.} El \textit{worker} asíncrono en
        proceso (sección~\ref{sec:orquestacion}) es suficiente para el \ac{MVP}, pero mantiene los
        trabajos en memoria: una caída del proceso los pierde y no admite más de una instancia. Una
        cola durable respaldada por Redis permitiría persistir los trabajos en curso, ejecutar
        varias réplicas del \textit{worker} en paralelo y, sobre esa base, incorporar mecanismos
        como la caducidad de las revisiones pendientes que hoy no tienen plazo.

  \item \textbf{Reintento de trabajos consciente del progreso ya realizado.} Actualmente, reintentar
        un trabajo fallido recupera el archivo de entrada original desde el almacenamiento de
        \textit{blobs} y repite el \textit{pipeline} completo desde cero, incluida una nueva llamada
        al proveedor de transcripción cuando el origen es audio. Sin embargo, cada etapa ya persiste
        su artefacto durable en \textit{blob storage}, de modo que la información para reanudar sin
        repetir trabajo ya existe; lo que falta es un reintento que la aproveche. La idea es sondear
        los artefactos en orden inverso de etapa, localizar el más avanzado que ya esté presente y
        reanudar desde ahí, evitando las llamadas externas redundantes y facturables.
\end{itemize}

\subsection{Extensiones del núcleo generativo}

Funcionalidades alineadas con los objetivos del sistema, pospuestas por falta de tiempo o por
requerir datos de producción para validarse.

\begin{itemize}
  \item \textbf{Arneses de evaluación para los componentes de recuperación aún sin cubrir}: el
        extractor de palabras clave, el generador \textit{multi-query} y el \textit{reranker}.
        Los tres son agentes ligeros de propósito único, análogos al de \textit{grounding}, por
        lo que un arnés propio, es decir, un corpus etiquetado más un \textit{scorer} determinista,
        encaja de forma natural en el marco existente. El más difícil de evaluar es el generador
        \textit{multi-query}: al pedirle reformulaciones libres de la consulta, la salida esperada
        no es única, lo que sugiere explorar señales como la \acs{NLI} —aplazada como señal de
        confianza (sección~\ref{sec:confidence-scoring}) pero aún sin analizar en este contexto—
        para juzgar si cada variante preserva la intención original.

  \item \textbf{Recuperación agéntica.} Un agente autónomo con herramientas de consulta como
        alternativa al \textit{pipeline} \ac{RAG}, que recupere nodos por razonamiento iterativo en
        lugar de por similitud vectorial.

  \item \textbf{Señal de confianza basada en \acs{NLI}.} Complementar las heurísticas con una
        segunda señal, ortogonal, que verifique si la descripción de un nodo se deduce de su cita
        fuente; se aplaza por requerir un modelo local con latencia aceptable y por quedar
        pendiente confirmar que aporta información realmente ortogonal a las heurísticas.

  \item \textbf{Detección y anonimización de PII.} Identificar y ocultar información personal en
        las transcripciones antes de procesarlas, para escenarios de mayor exposición que el
        despliegue interno previsto.

  \item \textbf{Captura de tareas sin responsable.} La identificación descarta hoy los ítems de
        acción sin un responsable nombrado; capturarlos con confianza reducida se pospone hasta
        contar con evidencia de su necesidad en transcripciones reales.
\end{itemize}

\subsection{Profundización del grafo de conocimiento}

Trabajo centrado en el modelo de datos, su ciclo de vida y su almacenamiento, más que en los
componentes de \ac{IAG}.

\begin{itemize}
  \item \textbf{Archivado de nodos (borrado lógico).} La base de conocimiento sigue hoy un modelo
        de sólo-adición: un nodo aprobado no puede retirarse salvo por la re-ingesta forzada, que lo
        elimina físicamente. Un estado \texttt{ARCHIVED} permitiría un borrado lógico reversible
        —excluir el nodo de la búsqueda y la navegación, y retirar su \textit{embedding} del almacén
        vectorial, sin perder su historial de auditoría—, de modo que la re-ingesta forzada
        archivara los nodos del trabajo anterior en lugar de destruirlos.

  \item \textbf{Eje de ciclo de vida de los nodos.} El modelo actual distingue el tipo y el estado
        de un nodo (vigente, supersedido, enmendado), pero no si el asunto que representa sigue
        \emph{abierto} o ya se ha \emph{cerrado}. Un tercer eje de ciclo de vida —una pregunta
        abierta que pasa a \texttt{RESOLVED}, un riesgo a \texttt{MITIGATED}, una decisión a
        \texttt{ENACTED} cuando una relación entrante la completa— enriquecería tanto las consultas
        al grafo como el filtrado de candidatos en la recuperación, evitando arrastrar a la
        resolución nodos ya zanjados. Su despliegue está previsto por fases, e implicaría ampliar el
        vocabulario de relaciones (por ejemplo, una arista \texttt{implements} de un ítem de acción
        hacia la decisión que materializa).

  \item \textbf{Revisión humana de las relaciones inferidas.} Una fase de revisión para las aristas,
        análoga a la que hoy existe para los nodos, que hoy se escriben sin aprobación explícita.

  \item \textbf{\textit{Backend} de grafo nativo.} La base de conocimiento se almacena hoy en
        PostgreSQL; un \textit{backend} de grafo nativo como \href{https://neo4j.com}{Neo4j}, tras
        la abstracción de almacén ya existente, habilitaría travesías y consultas sobre el grafo más
        expresivas que las que permite el modelo relacional actual.

  \item \textbf{Validación de integridad del grafo.} Detección de nodos duplicados, comprobación de
        integridad de las relaciones y guardas contra referencias circulares.

  \item \textbf{Umbrales del \textit{gate} por tipo de relación.} El \textit{gate} de resolución
        aplica hoy un umbral uniforme por tipo de nodo origen; discriminar por tipo de relación
        —una vez el corpus sea lo bastante grande— permitiría exigir más a las relaciones de
        frontera más nítida y menos a las ambiguas, en línea con la limitación de la coincidencia
        exacta de terna señalada en la sección~\ref{sec:analisis-errores}.
\end{itemize}

\subsection{Integración y nuevas fuentes de entrada}

Líneas que amplían los límites del sistema: qué formatos puede ingerir y con qué sistemas externos
se conecta.

\begin{itemize}
  \item \textbf{Entrada de vídeo y diarización de locutores.} Aceptar archivos de vídeo —con
        extracción de audio mediante ffmpeg— y segmentar la transcripción por locutor; ambas
        requieren muestras de producción para validar su calidad y un endurecimiento de las
        dependencias necesarias.

  \item \textbf{Inicialización de la base desde un corpus documental.} Sembrar la base de
        conocimiento a partir de documentación existente antes de la primera reunión, como
        herramienta complementaria de arranque. Sus fuentes naturales son las herramientas de
        gestión del conocimiento donde ya reside esa documentación, como
        \href{https://www.notion.so}{Notion} o \href{https://www.atlassian.com/software/confluence}{Confluence}.
\end{itemize}
